% ============================================================
% PRN232 - Section 2: Related Work (Vietnamese)
% ============================================================
\section{\bi{Công trình liên quan}{Related Work}}
\label{sec:related}

\bi{
Bài báo này định vị OmniCast trong bối cảnh ba hướng tiếp cận chính:
khung ứng dụng phía sau, giao thức truy cập dữ liệu, và hệ thống lịch
phát truyền thông.
}{
This paper positions OmniCast against three main lines of work:
back-end application frameworks, data-access protocols, and broadcast
scheduling systems.
}

\subsection{\bi{Khung ứng dụng phía sau}{Back-end Application Frameworks}}

\bi{
NestJS \cite{nestjs} chọn mô hình Dependency Injection dựa trên decorator
và mô hình module phân cấp, cho phép mỗi miền nghiệp vụ (kênh, chương
trình, người dùng) được đóng gói thành một module độc lập với vòng đời
riêng. So với Express thuần, framework này cung cấp sẵn lớp định tuyến
có phiên bản (\texttt{VersioningType.URI}), \texttt{ValidationPipe} với
chế độ \texttt{forbidNonWhitelisted}, \texttt{ExceptionFilter} và
\texttt{Interceptor} --- bốn thành phần mà bài toán phát sóng đòi hỏi
để mọi endpoint trả về cùng một định dạng phản hồi và mọi đầu vào
đều được kiểm tra nghiêm ngặt. Prisma \cite{prisma} đóng vai trò lớp
truy cập dữ liệu: \textit{schema} Prisma là nguồn sự thật duy nhất, và
\texttt{prisma migrate} sinh SQL di trúyền, nhờ vậy lịch sử thay đổi
cấu trúc bảng được kiểm soát phiên bản cùng mã nguồn.

Hai lựa chọn khác đáng so sánh là FastAPI (Python) và ASP.NET Core
(C\#). FastAPI tạo tài liệu OpenAPI tự động rất tốt nhưng hệ thống
đã có ngăn sách giới hạn ngôn ngữ; ASP.NET Core có hệ sinh thái
mature hơn nhưng cũng đòi hỏi đổi đội ngũ. NestJS cho phép tái sử
dụng ngôn ngữ TypeScript với cổng web Next.js \cite{nextjs}, nên một
kiểu dữ liệu hay một thay đổi hợp đồng API có thể được chia sẻ giữa
hai sản phẩm mà không cần bước chuyển đổi ngôn ngữ.
}{
NestJS \cite{nestjs} adopts a decorator-driven dependency-injection
model and a hierarchical module model, allowing each business domain
(channels, programmes, users) to be packaged as an independent module
with its own lifecycle. Compared with plain Express, the framework
provides out-of-the-box URI versioning (\texttt{VersioningType.URI}),
a \texttt{ValidationPipe} in \texttt{forbidNonWhitelisted} mode, an
\texttt{ExceptionFilter} and an \texttt{Interceptor} --- the four
components the broadcasting problem demands so that every endpoint
returns an identical response envelope and every input is validated
strictly. Prisma \cite{prisma} serves as the data-access layer: the
Prisma schema is the single source of truth and \texttt{prisma
migrate} generates version-controlled SQL, so schema drift is tracked
alongside the source code.

Two alternatives deserve comparison: FastAPI (Python) and ASP.NET
Core (C\#). FastAPI produces excellent OpenAPI documentation but the
project has a fixed language budget; ASP.NET Core has a more mature
ecosystem but also demands a language switch. NestJS allows the
TypeScript toolchain to be reused with the Next.js portal
\cite{nextjs}, so a shared data type or an API contract change can
propagate across both deliverables without a language boundary.
}

\subsection{\bi{Giao thức truy cập dữ liệu}{Data-Access Protocols}}

\bi{
Phần lớn hệ thống truyền hình truyền thống cung cấp lịch phát cho ứng
dụng di động thông qua OData \cite{odatapub}, một mở rộng của REST với
cú pháp truy vấn nhúng ($filter$, $orderby$, $expand$). OData rất phù
hợp khi client cần truy vấn linh hoạt trên một catalogue lớn mà không
muốn triển khai logic lọc phía server. OmniCast chọn hướng khác: các
truy vấn của EPG có hình dạng cố định (theo ngày, theo kênh, theo khung
giờ), vì vậy endpoint REST thuần với tham số truy vấn chuẩn hoá
(\cite{rest}) vừa đủ và đơn giản hơn cho cả hai client.

So với GraphQL \cite{graphql}, REST thuần giảm bớt chi phí phía client
vì không cần thêm một lớp runtime trên di động, đồng thời dễ cache
tầng HTTP và dễ kiểm thử bằng công cụ dòng lệnh. Đổi lại, hệ thống
chấp nhận việc phải thiết kế endpoint rõ ràng; bù lại, phần API được
tài liệu hoá hoàn toàn bằng OpenAPI 3 \cite{openapi} với decorator
\texttt{@ApiOperation}, \texttt{@ApiResponse}, \texttt{@ApiQuery} nên
không có tài liệu nào bị bỏ sót.
}{
Most legacy broadcasting systems expose their programme guide to
mobile applications through OData \cite{odatapub}, a REST extension with
embedded query syntax (\$filter, \$orderby, \$expand). OData is well
suited when clients need flexible queries over a large catalogue
without deploying server-side filtering logic. OmniCast takes the
opposite approach: EPG queries have a fixed shape (by day, by channel,
by time window), so plain REST endpoints with normalised query
parameters \cite{rest} are sufficient and simpler for both clients.

Compared with GraphQL \cite{graphql}, plain REST lowers the client-side
cost because no additional runtime layer is needed on mobile, and it is
easier to cache at the HTTP layer and to exercise from a command line.
The trade-off is that every endpoint must be designed explicitly; in
exchange, the entire API is documented with OpenAPI 3 \cite{openapi}
using \texttt{@ApiOperation}, \texttt{@ApiResponse} and
\texttt{@ApiQuery} decorators, so no endpoint is left undocumented.
}

\subsection{\bi{Lịch phát truyền thông và trợ lý AI}{Broadcast Scheduling and AI Assistance}}

\bi{
Hệ thống lịch phát thương mại (như EBU Radio-MEP và các bản ghi EPG của
FTA) bảo đảm tính toàn vẹn lịch bằng cách áp dụng thuật toán phát
hiện xung đột theo khoảng thời gian tương đối và bổ sung chương
trình lấn át khi truyền hình trực tiếp. Đây chính là mô hình mà
OmniCast tái hiện ở quy mô nhỏ hơn: khoảng thời gian của mỗi chương
trình là một đoạn $[t_0, t_1]$ trên trục thời gian của kênh, và hai
đoạn chồng nhau là không hợp lệ.

Về phía trợ lý biên tập, các hệ thống EPG thương mại gần đây bắt đầu
dùng mô hình ngôn ngữ lớn để sinh mô tả và thẻ nội dung
\cite{openai}. Điểm khác biệt của OmniCast là \emph{kiểm soát chất
lượng đầu ra}: mọi báo cáo do AI sinh ra đều được kiểm tra theo một
schema Zod \cite{zod} trước khi ghi vào cơ sở dữ liệu, kèm dấu vết
phiên bản prompt và phiên bản mô hình, nhờ vậy người quản trị có thể
so sánh chất lượng giữa các phiên bản prompt khác nhau.
}{
Commercial scheduling systems (such as EBU Radio-MEP and FTA EPG
feeds) guarantee schedule integrity by applying relative-time
conflict detection and substituting overrun programmes during live
transmission. This is precisely the model OmniCast reproduces at a
smaller scale: the slot of a programme is an interval $[t_0, t_1]$ on
the timeline of a channel, and two overlapping intervals are invalid.

On the editorial-assistance side, commercial EPG systems have recently
begun using large language models to generate descriptions and content
tags \cite{openai}. OmniCast's distinguishing feature is
\emph{output quality control}: every AI-generated report is validated
against a Zod schema \cite{zod} before being written to the database
and is stamped with the prompt version and the model version, so
administrators can compare quality across prompt revisions.
}

\subsection{\bi{Khoảng trống mà OmniCast lấp đầy}{Gap Addressed by OmniCast}}

\bi{
Tổng hợp lại, khoảng trống mà dự án lấp đầy là sự thiếu vắt một hệ
thống lịch phát truyền thông \emph{hoàn chỉnh và có thể kiểm thử tự
động} ở quy mô phòng học: vừa có cơ chế bảo đảm tính toàn vẹn lịch
phát, vừa có trải nghiệm người dùng hoàn chỉnh trên web và di động,
vừa có trợ lý AI có kiểm soát chất lượng, và toàn bộ được kiểm chứng
bằng 482 kiểm thử tự động trong quy trình tích hợp liên tục.
}{
In summary, the gap this project fills is the absence of a complete,
automatically verifiable broadcast-scheduling system at classroom
scale: schedule-integrity guarantees, a full user experience across
web and mobile, quality-controlled AI assistance, and the whole thing
validated by 482 automated tests inside a continuous-integration
pipeline.
}
